\pdfbookmark[1]{II. Untersuchung des gegenwärtigen Zustands der Professorenwitwenkasse zu Göttingen}{witwen-II}
\begin{center}
{\large [II.]\par}\medskip
\textit{Untersuchung des gegenwärtigen Zustands der Professorenwitwenkasse zu Göttingen.}\par\medskip
\textit{Vorwort.}\par
\end{center}
\noindent\hspace*{\parindent}In dem Votum, das ich am 8. Januar dieses Jahres in der Witwenkassenangelegenheit abgegeben habe, habe ich die wesentlichen Elemente der Methode angedeutet, die ich für die einzig geeignete halte, um zu einem so gründlichen Urteil zu gelangen, wie es die Natur des Gegenstands zulässt. Die dort bezeichneten, allerdings sehr langwierigen Rechnungen habe ich jetzt abgeschlossen; ihre Ergebnisse sind in der zweiten Abteilung dieser Denkschrift enthalten.

Da ich jedoch bei den meisten Mitgliedern des Kollegiums, dem diese Schrift vorgelegt wird, keine nähere Vertrautheit mit den Grundsätzen solcher Rechnungen voraussetzen kann, dachte ich, sie würden den Gegenstand gern auch noch von anderen Seiten und unter allgemeinverständlicheren Gesichtspunkten betrachtet sehen. Auch wenn sich auf diese Weise keine eigentlich \textit{präzisen}
\SourceLeaf{131}
Ergebnisse gewinnen lassen, sondern nur allgemeine Überschläge und Anhaltspunkte, ist es doch wichtig, diese im Einklang mit den Ergebnissen einer strengeren Methode zu finden; jedenfalls wird dadurch alles in ein helleres Licht gerückt. Außerdem hängen diese Erläuterungen eng mit der Überprüfung eines entscheidenden Punkts des bestehenden Regulativs zusammen. Diese Prüfung halte ich für unumgänglich notwendig, und dabei darf nichts im Dunkeln bleiben. Ich habe daher diese Entwicklungen im ersten Teil dieses Aufsatzes so ausführlich und, wie ich hoffe, so klar dargestellt, dass man ihnen leicht folgen kann.

